\section{Discussion}\label{sec:discussion}

PolyEdge represents a practical attempt to build autonomous AI for prediction market trading with hard safety boundaries. The system integrates three technical contributions---bounded AGI autonomy, evolutionary strategy composition, and dual-debate decision validation---into a single deployment pipeline. This section discusses the implications, limitations, and future directions of this work.

\subsection{Integration of Contributions}

The three contributions operate as complementary safety layers. The bounded autonomy framework (Section~\ref{sec:autonomy}) creates a sandbox in which experiments can run without risking production capital. Within that sandbox, the evolutionary strategy composition system (Section~\ref{sec:evolution}) generates novel strategies through mutation and crossover. The dual-debate validation system (Section~\ref{sec:debate}) then adversarially scrutinizes proposed trades before live deployment. Each layer reduces the probability of catastrophic errors: autonomy prevents runaway experiments, evolution constrains strategies to domain-specific chromosomes, and debate catches logical errors that numerical risk checks might miss.

This layered architecture mirrors the defense-in-depth approach advocated in safety-critical engineering \citep{nygard2018release}. No single safety mechanism is sufficient; the combination creates redundancy.

\subsection{Limitations}

We acknowledge several important limitations:

\textbf{Limited statistical power.} The 162 trades in the current dataset are insufficient to establish statistical significance for performance claims. The system is operational, but claims about its efficacy remain descriptive rather than inferential.

\textbf{Single deployment.} All data comes from one instance deploying to Polymarket and Kalshi.

\textbf{Rudimentary fitness.} The composite fitness function is linear, while true multi-objective optimization would require non-linear interaction, Pareto fronts, and regime-dependent weights.

\textbf{Static LLM evaluation.} The Judge agent in the debate system uses a fixed prompt and does not learn from historical debate outcomes. Future work could train the Judge on the TradeAttempt ledger.

\textbf{Hard-coded risk profiles.} The four risk profiles (safe, normal, aggressive, extreme) are static. Adaptive risk profiling that responds to market regime could improve capital efficiency.

\textbf{Negative P\&L.} The current deployment shows negative returns. While the bounded autonomy framework prevents catastrophic losses, it does not guarantee profitability.

\subsection{Future Work}

\textbf{Multi-platform expansion.} The current system targets Polymarket and Kalshi. Expansion to other platforms would diversify data sources.

\textbf{On-chain verification.} All trade settlement could be verified on-chain, providing immutable audit trails.

\textbf{Community governance.} The genome repository could become a public marketplace.

\textbf{Adversarial robustness.} Red-teaming the debate system with deliberately deceptive arguments.

\textbf{On-device LLM inference.} Running the debate Judge on-device would reduce latency and API costs.

\textbf{Meta-learning across strategies.} A meta-strategy could learn which strategies perform best in which regimes.

\textbf{Probabilistic safety guarantees.} From deterministic safety gates to probabilistic bounds on failure rates.

\textbf{Cross-institutional deployment.} Multiple instances sharing anonymized experiment results.

\subsection{Broader Implications}

Autonomous AI in financial markets is not a hypothetical future scenario---it is already being deployed. The question is not whether autonomous trading is coming, but whether it will be bounded. Our bounded autonomy framework demonstrates that hard safety boundaries are technically feasible: deterministic gates enforced in code, budget caps with automatic fallback, and experiment isolation. These mechanisms do not solve the full AGI safety problem, but they provide a pragmatic first step for narrow-domain autonomous systems.

The open-sourcing of the PolyEdge codebase, data, and experimental results contributes to the community goal of safe autonomous AI. By sharing both successes and failures (including the current negative P\&L), we hope to encourage a culture of transparency in autonomous trading research.
